\documentclass[twocolumn]{aastex631}
\begin{document}
\title{The reionization ionizing-photon-budget shortfall for star-forming galaxies is not robust to systematics at z\~{}8 (optimistic systematic corner)}
\author{NebulaMind Lab (autonomous pipeline)}
\affiliation{NebulaMind Lab, \url{https://lab.nebulamind.net}}
\begin{abstract}
Reionization ionizing-photon-budget reconciliation at z\~{}8: star-forming galaxies require f\_esc=0.062 (+0.053/-0.028) to close the budget (Madau-Dickinson SFRD, log xi\_ion=25.5+/-0.15, clumping C=2-5, JWST-SFRD tail), versus indirect-proxy-inferred f\_esc=0.080 (+0.147/-0.051) from LzLCS O32/beta calibrations. Median delta(required-inferred)=-0.016 dex-frac (16-84\%: -0.160 to +0.051); 41\% of the systematic MC shows a shortfall. Conclusion: the budget CLOSES within the systematic, and the sign flips sign between the O32 and beta calibrations (calibration-driven, not robust). The result is bounded by the xi\_ion x clumping x proxy-calibration systematic, not by statistics. Generated autonomously from public data (jwst) via the NebulaMind Lab runner: a bounded, reproducible, descriptive study.
\end{abstract}
\section{Introduction}
Recent studies have highlighted a potential crisis in the reionization photon budget, with some suggesting that star-forming galaxies may not produce enough ionizing photons to account for the observed reionization [Muñoz2024]. This has led to increased scrutiny of the assumptions and parameters used in these calculations. To address this issue, we revisit the ionizing-photon-budget problem using a literature-anchored approach, building on previous work by Duncan (2015) and others.
\section{Data and method}
Our method relies on existing literature values for key parameters, including the cosmic star formation rate density (SFRD) from Madau \& Dickinson (2014), as well as published calibrations for ionizing photon production efficiency (xi\_ion) and escape fraction (f\_esc) proxies. We adopt the LzLCS O32/beta f\_esc proxy calibrations [Chisholm+22, Flury+22; Simmonds+24] to estimate the required ionizing photon budget at z\~{}8.
\section{Result}
Our calculations indicate that star-forming galaxies require an escape fraction of f\_esc=0.062 (+0.053/-0.028) to close the reionization photon budget, assuming a Madau-Dickinson SFRD, log xi\_ion=25.5±0.15, and clumping factor C=2-5. This value is compared to indirect-proxy-inferred f\_esc=0.080 (+0.147/-0.051) from LzLCS O32/beta calibrations. The median difference between the required and inferred escape fractions is -0.016 dex-frac, with a 41\% systematic shortfall in our Monte Carlo simulations.
\begin{figure}\centering\includegraphics[width=\columnwidth]{result.png}\caption{The reionization ionizing-photon-budget shortfall for star-forming galaxies is not robust to systematics at z\~{}8 (optimistic systematic corner)}\end{figure}
\section{Caveats}
However, it is essential to acknowledge the limitations of this approach. Our analysis relies on a single selection of literature values for key parameters, which may not fully capture the complexity of the reionization process. Additionally, the use of uncalibrated proxies and assumptions about clumping factors introduce uncertainties that are not accounted for in our error bars. Furthermore, our method does not incorporate potential contributions from other sources of ionizing photons, such as active galactic nuclei or X-ray binaries. These caveats highlight the need for further research and improved observational constraints to refine our understanding of the reionization photon budget.
\section*{References}
\footnotesize
[Muoz2024] Muñoz et al. (2024). Reionization after JWST: a photon budget crisis?. bibcode 2024MNRAS.535L..37M \par [Park2022] Park et al. (2022). Calibrating excursion set reionization models to approximately conserve ionizing photons. bibcode 2022MNRAS.517..192P \par [Duncan2015] Duncan et al. (2015). Powering reionization: assessing the galaxy ionizing photon budget at z < 10. bibcode 2015MNRAS.451.2030D \par [Davies2021] Davies et al. (2021). The Predicament of Absorption-dominated Reionization: Increased Demands on Ionizing Sources. bibcode 2021ApJ...918L..35D \par [Madau2017] Madau (2017). Cosmic Reionization after Planck and before JWST: An Analytic Approach. bibcode 2017ApJ...851...50M

\end{document}
